\documentclass[11pt]{article}
\usepackage[T1]{fontenc}
\usepackage[utf8]{inputenc}
\usepackage{lmodern}
\usepackage[margin=1in]{geometry}
\usepackage{amsmath,amssymb,amsthm,booktabs}
\usepackage[hidelinks]{hyperref}
\newtheorem{theorem}{Theorem}
\newtheorem{lemma}{Lemma}
\newcommand{\JJ}{\mathcal J}
\newcommand{\FF}{\breve F}
\title{A finite counterexample to ideal-cardinality\protect\\ homomesy on constant-parameter fences}
\author{Aurantium}
\date{13 October 2026}
\begin{document}
\maketitle
\begin{abstract}
We give an explicit finite counterexample to Conjecture 5.4(b) of
Elizalde, Plante, Roby, and Sagan, \emph{Rowmotion on fences}.
For the 32-element fence $\FF(3^{11})$, the lower ideal
$\{5,6,7,8,12\}$ has rowmotion period 237, and the sum of ideal
cardinalities on its orbit is 3774. The orbit average is therefore
$1258/79$, rather than $16$. We supply the complete orbit, an independent
integer-arithmetic checker, and a Lean 4 formalization connecting the
paper's segment convention to the orbit calculation. The formalization
uses \texttt{native\_decide}; its trust assumptions include
\texttt{Lean.ofReduceBool} and the correctness of the Lean compiler.
No claim of priority is made.
\end{abstract}

\section{The conjecture and the precise fence}
For a finite poset $P$, write $\JJ(P)$ for its set of lower ideals.
Rowmotion on lower ideals is
\[
  \rho(I)=\mathord\downarrow\min(P\setminus I).
\]
A statistic $f$ is $c$-mesic if its average on every rowmotion orbit is
$c$. Conjecture 5.4(b) of \cite[p.~35]{eprs} states that, for a fence
$\FF(a^t)$ with constant segment parameter and an odd number $t$ of
segments, the statistic $I\mapsto |I|$ is $|P|/2$-mesic.
Segments here are maximal chains. The parameter of a segment is one plus
the number of its \emph{unshared} vertices, not uniformly its number of edges;
a shared vertex belongs to two adjacent segments \cite[p.~17]{eprs}.

Let $P$ have vertices $1,\ldots,32$. Its cover relations are given by the
following chains, each written from smaller to larger in the poset:
\[
\begin{array}{lll}
1<2<3, & 6<5<4<3, & 6<7<8<9,\\
12<11<10<9, & 12<13<14<15, & 18<17<16<15,\\
18<19<20<21, & 24<23<22<21, & 24<25<26<27,\\
30<29<28<27, & 30<31<32. &
\end{array}
\]
Equivalently, the underlying cover graph is the path $1,\ldots,32$.
For $1\le e\le31$, its edge from $e$ to $e+1$ points upward exactly
when $\lfloor e/3\rfloor$ is even. The partial order is the reflexive
transitive closure of these oriented edges.

The displayed chains are exactly the maximal chains: a chain can follow
the oriented path only until the next change of direction. Their shared
vertices are $3,6,9,\ldots,30$. Each endpoint segment has three vertices
and one shared vertex, whereas each of the nine internal segments has
four vertices and two shared vertices. Thus every segment has exactly
two unshared vertices and parameter 3. There are 11 segments and 32
vertices, so this is precisely $P=\FF(3^{11})$ in the convention of
\cite{eprs}. In particular, the number of segments is odd.

\section{The orbit certificate and the contradiction}
\begin{theorem}\label{thm:counterexample}
For the fence $P=\FF(3^{11})$ defined above, the lower ideal
$I_0=\{5,6,7,8,12\}$ has a rowmotion orbit $\mathcal O$ satisfying
\[
 |\mathcal O|=237,\qquad \sum_{I\in\mathcal O}|I|=3774.
\]
Its least positive return time is 237. Consequently Conjecture 5.4(b)
of \cite{eprs} is false.
\end{theorem}

The complete finite certificate appears in Appendix~\ref{app:certificate}.
An eight-digit hexadecimal integer $m_k$ encodes
\[
 I_k=\{v\in\{1,\ldots,32\}:\text{bit }v-1\text{ of }m_k\text{ is }1\}.
\]
The accompanying \texttt{data/orbit.csv} also lists the vertices,
cardinality, and maximal antichain of every ideal. The following
elementary certificate criterion explains exactly what is checked.

\begin{lemma}\label{lem:certificate}
Let $L$ be a positive integer. Suppose $I_0,\ldots,I_{L-1}$ are distinct lower ideals and
$\rho(I_k)=I_{(k+1)\bmod L}$ for every $0\le k<L$.
Then the forward orbit of $I_0$ is exactly this set of $L$ ideals,
and its least positive return time is $L$.
\end{lemma}
\begin{proof}
Induction gives $\rho^n(I_0)=I_{n\bmod L}$ for every $n\ge0$.
Thus every iterate is in the list, and every listed ideal occurs as an
iterate. Closure gives a return at $L$; distinctness excludes a return
at any $0<n<L$.
\end{proof}

\begin{proof}[Proof of Theorem~\ref{thm:counterexample}]
The principal lower sets of $5$, $6$, $7$, $8$, and $12$ are contained
in $I_0$, so $I_0$ is a lower ideal. For an explicit verification of the
certificate, define $D(x)=\{y:y\le_P x\}$. The order can be evaluated
without a closure algorithm: for $x<y$ numerically, $x\le_P y$ exactly
when every edge between them points right, and $y\le_P x$ exactly when
every such edge points left. Together with reflexivity this determines
all principal lower sets.

For each of the 237 encoded sets, verify
\[
 x\in I_k\ \Longrightarrow\ D(x)\subseteq I_k.
\]
Then compute, with indices taken modulo 237,
\[
 M_k=\{x\notin I_k:D(x)\setminus\{x\}\subseteq I_k\},
 \qquad I_{k+1}=\bigcup_{x\in M_k}D(x).
\]
These formulas are respectively the lower-ideal condition, the minimal
elements of the complement, and the definition of rowmotion. The
supplied checker verifies all 237 equalities, including the last-to-first
equality, distinctness of all 237 masks, and the cardinality sum 3774.
The Lean development verifies the same finite conditions on actual
finite sets and proves the conclusion of Lemma~\ref{lem:certificate}.
All arithmetic is integral or rational.

It follows that the complete orbit has mean
\[
 \frac{3774}{237}=\frac{1258}{79}=16-\frac6{79}\ne16=\frac{|P|}{2}.
\]
Equivalently, the conjectured sum is $237\cdot16=3792$, exceeding the
actual sum by 18. This single orbit satisfies the conjecture's hypotheses
and violates its conclusion.
\end{proof}

\section{Reproduction and formal verification}
From the root of the accompanying package, run
\begin{verbatim}
python checker/check.py
\end{verbatim}
The checker uses only the Python standard library. It reads the supplied
CSV instead of regenerating it, constructs the order directly from the
path orientation, and checks the complete orbit and its agreement with
the Lean orbit literals. The certificate is also fully printed here so
that the finite data can be checked independently of the distribution.

The Lean development fixes Lean 4.19.0 and mathlib at commit
\begin{center}
\small\texttt{c44e0c8ee63ca166450922a373c7409c5d26b00b}.
\end{center}
All transitive dependencies are pinned in the lockfile. Portable instructions, the
actual build logs, and the axiom audit are included. The four modules
have the following roles:
\begin{itemize}
\item \texttt{Poset}: the 32-element partial order; all maximal chains;
      distinct segments, shared vertices, and the parameter convention.
\item \texttt{Rowmotion}: genuine lower ideals and the identity
      $\rho(I)=\mathord\downarrow\min(P\setminus I)$.
\item \texttt{Orbit}: all 237 lower ideals, all transitions including
      closure, distinctness, and the sum 3774.
\item \texttt{Main}: every iterate, completeness of the orbit, its
      least period and mean, and the failure of half-cardinality homomesy.
\end{itemize}
The theorem \texttt{W04.paper\_fence\_counterexample} combines the
segment convention with the dynamical failure. The order on the vertices
is the explicit fence relation, rather than the numerical order on
\texttt{Fin 32}.

\paragraph{Trust assumptions.}
The finite orbit certificates use \texttt{native\_decide}.
The final counterexample theorem's printed axiom dependencies are
\begin{center}
\texttt{propext, Classical.choice, Lean.ofReduceBool, Quot.sound}.
\end{center}
In particular, \texttt{Lean.ofReduceBool} adds trust in native evaluation
and the correctness of the Lean compiler. This is not a completed
kernel-reduction-only proof. The final sources contain no
\texttt{sorry}, \texttt{admit}, or custom unproved mathematical axioms;
the final theorem dependencies contain no \texttt{sorryAx}. The CSV and
Python checker are not introduced as Lean axioms. Successful compilation
is a machine verification result, not a claim of external human review.

\section{Scope and literature boundary}
The conclusion concerns Conjecture 5.4(b), about ideal-cardinality
homomesy; it makes no assertion against Conjecture 5.4(a), about
antichain-cardinality homometry. No minimality among counterexamples is
claimed.

Mertin and Poznanovi\'c \cite{mp} restate the target as Conjecture 12
and study equivalent or stronger forms. Their three-segment result
(Theorem 19) and their result for dihedral orbits (Theorem 26) do not
assert the required average on each eleven-segment rowmotion orbit.
The present counterexample does not depend on their identities.
The literature check did not identify an earlier instance of this
counterexample in the materials inspected; this is not an exhaustive
priority determination. The official record and abstract of Mertin's
2025 dissertation \cite{mertin} were accessible, but its full text could
not be retrieved because the download returned HTTP 403. Its contents
have not been ruled out as a prior source of the result. We therefore
make no first-discovery claim.

\paragraph{AI assistance.}
This work was prepared with AI assistance. No independent human
peer-review claim is made.

\begin{thebibliography}{9}
\bibitem{eprs}
S. Elizalde, M. Plante, T. Roby, and B. E. Sagan,
\emph{Rowmotion on fences}, Algebraic Combinatorics \textbf{6}(1)
(2023), 17--36. \href{https://doi.org/10.5802/alco.256}{doi:10.5802/alco.256}.

\bibitem{mp}
A. Mertin and S. Poznanovi\'c,
\emph{Toggleability Spaces of Fences}, Electronic Journal of
Combinatorics \textbf{31}(4) (2024), P4.46.
\href{https://doi.org/10.37236/13085}{doi:10.37236/13085}.

\bibitem{mertin}
A. Mertin, \emph{Homomesies and Toggleability Spaces}, Ph.D. dissertation,
Clemson University, August 2025.
\href{https://open.clemson.edu/all_dissertations/4090/}{Official record 4090}.
Record and abstract inspected; full text not retrieved.
\end{thebibliography}

\appendix
\section{The complete orbit in hexadecimal}\label{app:certificate}
Read the entries from left to right and then top to bottom, starting at
$k=0$. There are 237 entries, with eight entries per full line. Bit $v-1$
encodes vertex $v$. Thus the first entry encodes $\{5,6,7,8,12\}$.
The last entry maps back to the first under rowmotion. The sum of the
numbers of set bits is 3774.
\begin{center}
\small\ttfamily
\begin{tabular}{@{}llllllll@{}}
000008f0 & 20821c39 & 71c73e63 & fbef80f0 & 3ffe0839 & 60031c63 & f087bef0 & 39cffff9 \\
63e00003 & f0020820 & 38871c71 & 61cfbefb & f3e3ffff & 38060000 & 608f0821 & f1c39c73 \\
3be63ef8 & 7f8f0fe1 & e0c39833 & 31e63c78 & 7b8f0ee1 & ffc39ff3 & 00e63838 & 218f0c61 \\
73c39ef3 & f8e63ff8 & 018f0001 & 23c38823 & 70e61c70 & f98f3ef9 & 03c38fe3 & 20e61830 \\
718f3c79 & fbc38ee3 & 3fe61ff0 & 600f3839 & f0838c63 & 39c61ef0 & 63ef3ff9 & f0ff8003 \\
39800820 & 63c21c71 & f0e73efb & 398f8fff & 63c01800 & f0e23c21 & 39870e73 & 63cf98f8 \\
f0e03c01 & 39820e23 & 63c71870 & f0efbcf9 & 39fffe03 & 63800020 & f0c20871 & 39e71cfb \\
638fbe3f & f0c3f860 & 39e60cf1 & 638f1e3b & f0c3b87f & 39e7fce0 & 638e0e31 & f0c3187b \\
39e7bcff & 638ffe00 & f0c00021 & 39e20873 & 63871cf8 & f0cfbe01 & 39e3f823 & 63860c70 \\
f0cf1ef9 & 39e3bfe3 & 6387f830 & f0ce0c79 & 39e31ee3 & 6387bff0 & f0cff839 & 39e00c63 \\
63821ef0 & f0c73ff9 & 39ef8003 & 63fe0820 & f0031c71 & 3887befb & 61cfffff & f3e00000 \\
38020821 & 60871c73 & f1cfbef8 & 3be3ffe1 & 7f860033 & e0cf0878 & 31e39ce1 & 7b863e33 \\
ffcf0078 & 00e388e1 & 21861c33 & 73cf3e78 & f8e380e1 & 01860833 & 23cf1c78 & 70e3bee1 \\
f987fff3 & 03ce0038 & 20e30861 & 71879cf3 & fbce3e38 & 3fe30061 & 600788f3 & f08e1c38 \\
39c33e61 & 63e780f3 & f00e0838 & 38831c61 & 61c7bef3 & f3effff8 & 38fe0001 & 61830823 \\
f3c79c70 & 38ee3ef9 & 61ff0fe3 & f3839830 & 38c63c79 & 61ef0ee3 & f3ff9ff0 & 38003839 \\
60820c63 & f1c71ef0 & 3befbff9 & 7ffff803 & e0000c20 & 30821e71 & 79c738fb & e3ef8c3f \\
30fe1e60 & 798338f1 & e3c78c3b & 30ee1e7f & 79ff38e0 & e3838c31 & 30c61e7b & 79ef38ff \\
e3ff8c00 & 30001e21 & 78823873 & e1c70cf8 & 33ef9e01 & 78fe3823 & e1830c70 & 33c79ef9 \\
78ee3fe3 & e1ff0030 & 33838879 & 78c61ce3 & e1ef3e30 & 33ff8079 & 780008e3 & e0821c30 \\
31c73e79 & 7bef80e3 & fffe0830 & 00031c79 & 2087bee3 & 71cffff0 & fbe00039 & 3f820863 \\
60c71cf0 & f1efbe39 & 3bfff863 & 7f800cf0 & e0c21e39 & 31e73863 & 7b8f8cf0 & ffc01e39 \\
00e23863 & 21870cf0 & 73cf9e39 & f8e03863 & 01820cf0 & 23c71e39 & 70efb863 & f9fffcf0 \\
03800e39 & 20c21863 & 71e73cf0 & fb8f8e39 & 3fc01863 & 60e23cf0 & f1870e39 & 3bcf9863 \\
7fe03cf0 & e0020e39 & 30871863 & 79cfbcf0 & e3e3fe39 & 30060063 & 788f08f0 & e1c39c39 \\
33e63e63 & 780f00f0 & e0838839 & 31c61c63 & 7bef3ef0 & ffff8ff9 & 00001803 & 20823c20 \\
71c70e71 & fbef98fb & 3ffe3c3f & 60030e60 & f08798f1 & 39ce3c3b & 63e30e7f & f00798e0 \\
388e3c31 & 61c30e7b & f3e798ff & 380e3c00 & 60830e21 & f1c79873 & 3bee3cf8 & 7fff0e01 \\
e0039823 & 30863c70 & 79cf0ef9 & e3e39fe3 & 30063830 & 788f0c79 & e1c39ee3 & 33e63ff0 \\
780f0039 & e0838863 & 31c61cf0 & 7bef3e39 & ffff8063 \\
\end{tabular}
\end{center}
\end{document}
